\section{Neo-Vedanta, Evil Spirits, and Don Quixote}

I attended a small gathering of Neo-vedantists last week. (A neo-vedantist is a Westerner who tries to extract the metaphysics of the Vedanta apart from their connections to the Vedas.) There was the typical talk about the higher self, which transcends space and time. Not bad, but not quite the same as actually transcending space and time.

Then there was talk about the mind rambling on at random. Again OK, but practices to master the mind was missing. The most experience woman asserted that as a brute fact, but could offer no explanation. That is because the neo-vedantists have no doctrine of the ``Fall'', so there is no understanding of how man in a higher state could fall into a lower state. Without that understanding, it seems more difficult to reverse the process. Also, in that system there is no way to understand Redemption.

Then the next topic was the world as illusion, which easily leads to nihilism. If everything is illusion, as these neo-vedantists claimed, then life is pointless. IF only consciousness is real, but why make the effort to remain fully conscious of a meaningless world for all eternity? The sounds more like hell than bliss!

Of course, every real spiritual path requires moral discipline in order to purify the mind. This is opposite to nihilism. Of course, the false paths are all illusory, but efforts must be made to follow the real.

\paragraph{Horizontal Memory and Universalism}


The next topic questioned the reality of time: The past no longer is, the future is not yet, the now disappears. Although such assertions are common in certain circles, does such a pseudo-spirituality make sense? If Gnosis is remembering, then what exactly are we remembering?

Hermetic teaching refers to vertical and horizontal memory, as described by \textbf{Valentin Tomberg}. Vertical memory brings what is above into present consciousness. This is more Platonic, since vertical memory is knowledge of the forms or prototypes. Horizontal memory, on the other hand, makes the past present.

If the past is ``not real'', then how can it be made present? If I think of a unicorn, I am visualizing something unreal. But when Achilles has a vision of his dead mother, is that at all analogous? Achilles is making the past present, whereas my vision is totally illusory.

One of the goals of the modern world is to render the past unreal. Human beings derive their being horizontally as well as vertically. Vertically, we share in a universal human nature, e.g., we instantiate certain prototypes. But that does not exhaust being, since we are born into specific familial, historical and social situations. I used to have to drop off my teenage son a block away from his friends' houses, as if his friends would be fooled into thinking he spontaneously arose and was not the fruit of sexual reproduction.

That is the common delusion today. We emphasis the universal but neglect the specific. We pretend our genetic inheritance is an unfortunate accident. We suppress the mythological-spiritual sources that define us. We try to become the abstract I of the neo-vedantists, that observes but does not participate.

\paragraph{Other Minds}


New atheists scoff at the idea of the world being populated by spirits. A spirit is an I, with conscious experience, a rational nature, and some level of moral agency. Hence, there are presumably several billion spirits inhabiting the earth. These human spirits are tied to bodies, and their everyday commonness shrouds a real mystery.

That is because science has no way to prove the existence of spirits. The only plausible answer is that consciousness is an illusion, an epiphenomenon of biochemical processes in the brain. Alan Turing proposed a test to distinguish a computer simulation of a person from a real person. If his test works, that just demonstrates that science cannot distinguish between a machine and a human being.

Philosophers debate the ``problem of other minds'' endlessly, without conclusion. Yet they don't doubt their existence, or else they would be logically compelled to embrace solipsism. I still recall a thought I had when I was a very young boy driving back from Grandma's house with my family. I was looking out at all the other cars, the houses, the people, and so on. I assumed that only I existed, that the only significant world event was dinner at Grandma's. If so, then who were all those other people and what were they doing? I considered that God was spontaneously creating the surrounding world simply to entertain me. That must have baffled that boy for a while, since I still more or less think that way. Only grudgingly have I learned to accept the existence of other people. Isn't it obvious? A solipsist would not bring so much pain and disappointment on himself, would he?

\paragraph{Evil Spirits}


I won't repeat the arguments that the human spirit can exist apart from the physical body, but simply stipulate it as a fact. If so, then there is no reason to doubt the possibility of the existence of other disembodied spirits. That is the logical argument, but there is also the empirical experience of entire history of the human race that testifies to their existence.

Those who have undergone Hermetic training, which involves careful observation and cataloging of the inner states of consciousness, eventually come to the realization that their thoughts are not necessarily their own. They begin to realize that it is as though another being is thinking those thoughts for them. Passive ideas have no power. \textbf{Giovanni Gentile} emphasizes this:

\begin{quotex}
The primary causal agent is always an idea become person, with a will that pursues determinate ends---a cognizant will that has a program to realize, a concrete thought, effective in history.
\end{quotex}


If that is true about human persons, it is a fortiori true about extra-human agents. The ideas that rattle around in our heads always have a person pushing it for ``determinate ends''. This may sound odd to modern men, but that is simply because they never bother to consider the actual sources of their thoughts.

\textbf{Sergei Bulgakov} claims that one's guardian angel is a prototype for the human. This is consistent with Rene Guenon's claim that the angels are higher states of the being. So we can become aware of angelic intelligences by reaching these higher states. Opposed to them are, let's call them ``devils''. There is actually a certain comfort in them, since devils believe in God and seem to have a compact with God limiting their effects. Furthermore, they are at least interested in us.

Then there may be a large number of beings that are totally indifferent to us. A horrifying death to me has always been the fear of being eaten by an animal, such as an alligator. That is because of the animal's complete indifference to us as a human being. At least a murderer knows he is killing another human.

Although the real Hermetic goal is the achievement of higher states, there are many more inferior schools that try to cultivate relationships with these indifferent beings. Some they try to control, others they need to supplicate. Still others claim to bring new teachings to humanity; they are not angels, so the results are mixed. There is today a proliferation of psychic channeling.

An interesting case is \textbf{H. P. Lovecraft} (since I recently acquired Barnes and Noble's attractive edition of the Cthulhu tales). Although Lovecraft was a professed atheist and materialist,
\textbf{Kenneth Grant}, in \emph{The Magical Revival}, has a different take on it. He wrote:

\begin{quotex}
Understandably terrified of crossing the Abyss, Lovecraft forever recoiled on the brink, and spent his life in a vain attempt to deny the potent Entities that moved him. Little wonder that the tales he wrote are among the most hideous and powerful ever penned.
\end{quotex}


Grant himself does not deny such Entities. He elaborates on their danger:

\begin{quotex}
Not only can the disembodied spirit of dead or sleeping people impersonate spirits and work evil by such means, but---which is infinitely more dangerous, extracosmic entities can masquerade as spirits and, if they are not banished before they can gain a foothold in the consciousness of he sic who invoked them, obsession follows. ... Lovecraft is the authority of the devastation they leave in their wake when they are let loose upon the earth.
\end{quotex}


It is curious that only some occultists are aware of the spiritual battle going on in the world. The proliferation of these extracosmic entities is the manifestation of the ``fissure in the great wall'' that Guenon described:

\begin{quotex}
the `fissures' occur only at the base, and therefore in the actual protective wall itself, and the inferior forces that make their way in through them meet with a much reduced resistance because under such conditions no power of a superior order can intervene in order to oppose them effectively. Thus the world is exposed defenceless to all the attacks of its enemies, the more so because, the present-day mentality being what it is, the dangers that threaten it are wholly unperceived.
\end{quotex}


Everyone else is like Don Quixote fighting imaginary battles in a world they no longer understand.


\flrightit{Posted on 2017-05-04 by Cologero}

\begin{center}* * *\end{center}

\begin{footnotesize}
\begin{sffamily}
\texttt{Osbourne on 2017-05-04 at 06:26 said:}

Shankara outlines in his Vivekachudamani (Jewel of Discernment) the necessary preparatory work for reintegration. The prerequisites to the work necessary are the antitheses to what these self-styled neo-vedantin blabbermouths spew out. Foremost is the discrimination between the real and unreal. The unreal subsumes under itself all the preconceived intellectual formulations arbitrarily constructed without achieving the corresponding state. The end result is far from guaranteed. It's not to be partaken of democratically. It's clear that these neo-vedantins (better inverted-vedantins) understand nonduality as a form of monism where everything is the same and equal under the common denominator of being an illusion. No discernment is necessary, just repeating the same illusion mantra like pacman. All these inverted spiritualities are in full conformity to the modern capitalistic zeitgeist. Zizek said aptly that in a CocaCola advertisement when they say ``That's It'', the ``It'' refers to a transcendent surplus quality. However, when you leave a coke bottle on the beach for 40 minutes, you find that it tastes like the sh*t that it actually is. In a capitalistic/materialistic society, what is marketed as the transcendental is actually the excremental. From this angle, one wonders if it is not natural that one neo-vedantin guru called Adidam (quite a caricature character) was having his female bimbo disciples and wives of his male disciples defecating in the bed during the act of coitus under the pretext that this work was necessary to transcend the limitations of the individual non-existent ego. After all, everything is an illusion and no discernment is necessary.

This hilarious character has one pseudo-intellectual admirer in the form of one Ken Wilber who spews out such illuminating passages as ``Aristotle, Plotinus, Saint Augustin, Hegel, Schuon, Coomaraswamy, Guenon were all developmentalists''. The developmentalism here is in the evolutionistic sense. I guess I should be taking some crazy pills. These types are either morons that don't know how to read or such sinister individuals that are so deluded with their misconceptions that they project them to others. After writing many tomes to glorify his search for a theory of everything (mimicing the modern scientific paradigm), which ouevre he arbitrarily divides into periods, he ends up in the last period revealing that there is no such thing as a separate spiritual realm.

In the eternal night of the neo-vedantins, everything is black. One cannot but admire Guenon for his foresight and diligence in anticipating potential areas of distortion and corruption of the metaphysical doctrine. It was not for no reason that he was so critical of Vivekenanda, who by the way looks like a saint compared to these modern neo-vedantins. This is proof like no other that in matters of doctrine, no concession can be given.

\hfill

\texttt{Schoolboy on 2017-05-07 at 22:49 said:}

Excellent article! I have experience with not only neo-advaita but also Buddhists. While the advaitans fall in the nothing is real category I have decided that Buddhists fall in the liberal, democratic wing of the American left. I've come to despise Buddhists as phonies and advaitans as being unable or unwilling to answer simple questions concerning spiritual life. I read Tomberg years ago before most people knew his real name, I found most of the Tarot explanations good, some great, others overdone. Wilber loved Adida and to this day defends him and his sexual excesses in Fiji. The whole Buddhist/advaita thing is a diversion from the destruction of the West, hoping the whole thing will go away. Regarding the reality of spirits existing outside of the corporeal existence, the whole idea of praying to saints for help, praying to ancestors for guidance has been around in all cultures in the world. With Buddhists/advaitans contacting spirits is passed', something Buddha would never consider or something Sankara would say is devoid of value, being illusory to begin with. Haven't read Guenon yet, will pick him up. Thank you, Michael

\hfill

\texttt{Mark Citadel on 2017-05-08 at 17:43 said:}

As always, Cologero, the profundity of your writing hits the mark perfectly. ``One of the goals of the modern world is to render the past unreal.'' but of course! The very reason the moderns describe prior ways of living with such terms as `primitive' and `backward', and even have gone so far as to turn Medieval into a kind of slur, is evidence of this wish to reduce the past to a kind of Gothic horror novel. When they attack the reactionary as wishing to return to a ``past that never existed'', they are engaging in precisely this kind of revisionism, all in order to affirm their numerous crimes.

\hfill

\texttt{THAT on 2017-06-05 at 21:06 said:}

Yes, it is funny how solipsistic certain pseudo-esoterists suddenly turn when non-human invisible beings of the intermediate world are concerned: then demons are ``only symbolical'' and so forth. The Islamic tradition is a particularly clear of the actual situation in this case, as it affirms the existence of a great variety of non-angelic intelligent beings classed as Jinn situated in close proximity to our human world, some of whom actively side with the satanic powers, while others may be servants of God or at least not hostile to humanity.

There are occultists who know about the war of which you speak and who consciously choose allegiance to the evil influences, which Guénon's ``counterinitiation'' essentially signifies. A publicly available resource that exemplifies how such a disturbing phenomenon might present itself in the most explicit form is the body of writing associated with the Order of Nine Angles underground cult, whose orientation deliberately aims at unleashing the ``dark gods'' of the abyss upon the world through so-called ``nexions'' constituted by voluntarily possessed initiates.

\hfill

\texttt{Fred on 2017-06-15 at 16:22 said:}

According to the writings of a small dead italian hermetic christian school the guardian angel is precisely the nous, the angelic intellect, buddhi.

P.S. the hermetic school is Erim's, a roman aristocratic of the past century (Guénon asked De Giorgio informations about him in a letter).

\hfill

\texttt{Cologero on 2017-06-18 at 07:10 said:}

Fred, I can't find anything about Erim. I've published most of the correspondence between Guido De Giorgio and Guenon ... which letter are you referring to?

\end{sffamily}
\end{footnotesize}

